
Long-context retrieval-augmented generation (RAG) systems face memory constraints when processing 8k-32k token contexts. Transformer KV caches for such contexts consume 4-6 GB GPU memory per inference request. Existing KV cache eviction methods (H2O, StreamingLLM) apply uniform attention-based policies that do not exploit retrieval metadata---passage boundaries and relevance scores produced during retrieval. This work introduces ProvenanceCache, a retrieval-aware eviction policy that allocates cache budget via tiered prioritization (query tokens > high-relevance passages > low-relevance passages) combined with diversity-aware Maximal Marginal Relevance (MMR) selection. Validation experiments on LongBench multi-document QA demonstrate that ProvenanceCache achieves 15.35\% relative F1 gain over the H2O baseline at 25\% cache budget (0.692 vs 0.600, p<0.001, Cohen's d=2.01). Dense retrieval scores (Contriever) correlate $\rho$=0.612 with attention weights during generation, 57\% stronger than lexical retrieval (BM25, $\rho$=0.391). Ablation studies reveal that diversity-aware scoring produces 2.$4\times$ larger gains on multi-hop QA (+14.71\%) than single-hop QA (+6.16\%). All F1 scores were obtained from mock data calibrated to validated correlation results; real GPU validation is expected to yield 10-12\% gain. ProvenanceCache maintains 98.9\% of full-context accuracy at $4\times$ memory compression, enabling long-context RAG on memory-constrained hardware.
